\documentclass{USN-BSc}

\usepackage[style=apa, backend=biber]{biblatex}
\addbibresource{thesis.bib}

% --- Fyll inn forsideinformasjon ---
\faculty{Fakultet for\dots.}
\subjectinfo{Emnenavn/emnekode}
\semester{Årstall og semester}
\candidate{[Navn/kandidatnummer]}
\title{[Tittel]}
\subtitle{[Undertittel]}

\begin{document}

\USNcover

% --- Innholdsfortegnelse ---
\tableofcontents
\clearpage

% =====================================================
\section{Innledning}

Denne malen er et forslag til hvordan du kan sette opp oppgaven. Får du
retningslinjer fra faglærer eller fra studiet følger du disse.

Disse retningslinjene gir generelle føringer for hvordan skriftlige oppgaver
kan struktureres. Samtidig kan de brukes som en mal for skriving av oppgaver,
der studenten erstatter teksten og overskriftene i disse retningslinjene med
sitt eget innhold.

% =====================================================
\section{Innhold og struktur}

Rekkefølgen på oppgavens innholdselementer skal ha en oppbygning som denne
(sammendrag og forord er kun aktuelt ved større oppgaver, og oppgavens lengde
og spørsmålsformulering avgjør hvilke elementer som skal være med):

\begin{itemize}
  \item[] Forside
  \item[] Evt.\ sammendrag
  \item[] Evt.\ forord
  \item[] Innholdsfortegnelse med sideangivelser
  \item[] Innledning
  \item[] Metode
  \item[] Teori
  \item[] Analyse
  \item[] Drøfting
  \item[] Konklusjon
  \item[] Litteraturliste
  \item[] Evt.\ vedlegg. Vedlegg nummereres, og vedlegg som har flere sider
          skal sidenummereres.
\end{itemize}

% =====================================================
\section{Overskrifter}

Overskrifter skal ha en logisk struktur etter hva som er mest viktig.
Hovedkapittelet markeres med Overskrift~1 som over dette avsnittet. Strukturen
skal gjøres synlig ved å bruke tall for inndeling av overskriftene, og ved å
bruke en skriftstørrelse og eventuelt skrifttype som skiller seg ut fra annen
tekst og reflekterer hva som er mest viktig. Overskriftene skal være
venstrejustert. Det skal ikke være punktum etter overskrift. I større oppgaver
skal hovedkapitlene starte på ny side. Det vil være naturlig å bruke
overskrift~1, 2 og eventuelt 3, som eksemplifisert i dette dokumentet.

Hvordan legge inn flere kapitler: bruk \verb|\section{...}| for Overskrift~1,
\verb|\subsection{...}| for Overskrift~2, og \verb|\subsubsection{...}| for
Overskrift~3. Overskrift~1 starter automatisk på ny side.

\subsection{Overskrift 2}

Overskrift~2 brukes til hovedpunkt innenfor et kapittel.

\subsubsection{Overskrift 3}

Overskrift~3 brukes til underpunkt innenfor hovedpunkter (overskrift~2).

% =====================================================
\section{Kildehenvisninger, sitering og litteraturliste}

Alle oppgaver skal ha en litteraturliste. Denne sorteres vanligvis alfabetisk
på forfatter.

\subsection{Kildekompasset}

Det finnes mange formelle systemer for å skrive en korrekt litteraturliste. Et
av de mest brukte systemene er APA. Kildekompasset har en utmerket
eksempelsamling i APA, se \url{https://kildekompasset.no/}. Andre systemer som
er mye brukt er Harvard, Vancouver og IEEE.

\subsection{Henvisninger til KI}

Dersom det er brukt KI i arbeidet med bacheloroppgave er det naturlig at det
skrives en mer omfattende tekst om hvordan KI er brukt og hvordan det har
påvirket arbeidet med oppgaven. Lag en egen del i oppgaven, for eksempel med
overskriften ``Bruk av kunstig intelligens'', eller legg det som fotnote.

Se mer informasjon om USN sine retningslinjer for bruk av KI her:
\url{https://www.usn.no/studier/studiekvalitetsuka/eksamen-fusk/retningslinjer-for-bruk-av-kunstig-intelligens-ki-ved-eksamen-og-studentoppgaver}

% =====================================================
\section{Etiske retningslinjer for skriftlige studentoppgaver}

Disse etiske retningslinjene gjelder for alle oppgaver som utføres av studenter
ved Universitetet i Sørøst-Norge. Spesielle regler for den enkelte oppgave kan
i tillegg være angitt fra fakultetet.

Studentoppgaver knyttet til praksisstudier eller kliniske studier vil som regel
inngå som en del av den ordinære og etablerte virksomheten på praksisstedet.
Enkelte oppgaver vil derimot kunne ha preg av å være forskning. Det vil si at
studentens arbeid er planlagt å omfatte innsamling av data og/eller inneholde
elementer av forsøk, hvor målet er utvikling av ny kunnskap. I disse tilfeller
må det vurderes om det bør søkes om tillatelse og prosjektplan legges fram for
Regional Etisk Komité (REK) for vurdering. Det er et ansvar for både student,
praksisstedets veileder og universitetet å nøye vurdere hvordan det enkelte
prosjekt skal forholde seg til dette. Et første skritt er å orientere seg på
heimesida til Etiske forskningskomitéer:
\url{https://www.forskningsetikk.no/} for å vurdere om det aktuelle arbeidet
faller innenfor mandatet til komiteene.

\subsection{Anonymisering av data}

Studenten har taushetsplikt. Studentens læresituasjon er underordnet generelle
bestemmelser om personvern. Alle person- og stedsdata som benyttes i en
undersøkelse, et prosjekt eller i en studentoppgave må anonymiseres slik at
ikke person eller sted kan identifiseres i den skriftlige
rapporten/besvarelsen. Anonymisering av opplysninger er studentens ansvar. For
å bidra til at personer og sted ikke kan identifiseres må studenten om
nødvendig avstå fra eller omskrive data. Manglende ivaretakelse av kravet om
anonymisering kan medføre at studentens arbeid blir vurdert til ikke bestått.

\subsection{Formell tillatelse/samtykke}

I de tilfeller hvor studentens arbeid omfatter selve praksisstedet og dets
ansatte, skal samtykke innhentes fra ledelsen og eventuelt den enkelte ansatte
som det skal innhentes data fra. Ledelsen for praksisstedet kan sette grenser
for studentens arbeid/datainnsamling.

\subsection{Offentliggjøring av oppgaver}

Universitetet kan benytte deler av godkjente oppgaver til undervisning, og kan
benytte oppgaver i forskningssammenheng. Publisering av hele oppgaver eller
bruk av studentarbeider til andre formål kan kun gjøres med særskilt tillatelse
fra vedkommende student(er). Student eller høgskole har muligheter for å
klausulere et skriftlig studentarbeid. Begrunnelsen vil i slike tilfeller være
at innholdet i rapporten kan skade en tredjepart om rapporten ble offentlig
tilgjengelig, selv ved forsøk på å anonymisere innholdet.

% =====================================================
\USNsection{Litteraturliste}

\nocite{*}
\printbibliography[heading=none]

\end{document}
